% practice_answers.tex -- Module 8 (From the Universe to the Atom): worked answers to the
% Practice set. One practiceans per practiceq in practice.tex, in the same order (25 answers).
% Local helpers (prefix "prac"), the same as in practice.tex (\providecommand: whichever
% file is read first defines them).
\providecommand\pracNuc[3]{\ensuremath{\mathord{%
  \setbox0=\hbox{$\scriptstyle#1$}\setbox2=\hbox{$\scriptstyle#2$}%
  \ifdim\wd0<\wd2 \setbox0=\hbox to\wd2{\hss\box0}\else\setbox2=\hbox to\wd0{\hss\box2}\fi
  {}^{\box0}_{\box2}\mathrm{#3}}}}
\providecommand\pracaq[1]{\ensuremath{\overline{\mathrm{#1}}}}

\begin{practiceanswers}

% ===================================================================== Section I
% ---- Q1 radiation to matter
\Needspace{9\baselineskip}
\begin{practiceans}
  \correct{C}
  In the first moments, the photons had huge energies. A photon with energy of at least
  $2mc^2$ can turn into a particle of mass $m$ and its antiparticle (pair production).
  Energy becomes mass, as $E = mc^2$ says. Mass--energy is conserved. As the Universe
  cooled, quarks joined into protons and neutrons, and later these made nuclei.
  \begin{itemize}
    \item \xmark\ A: atoms formed much later, about \num{380000} years after the Big Bang,
          when the Universe had cooled to about \qty{3000}{K}. Quarks, protons, neutrons and
          nuclei came first.
    \item \xmark\ B: the mass came from the photons' energy. Mass--energy was conserved.
    \item \xmark\ D: protons and neutrons formed in the first second. The first stars formed
          about 200 million years later. Stars fuse nuclei that already exist.
  \end{itemize}
\end{practiceans}

% ---- Q2 Curie-Joliot
\Needspace{9\baselineskip}
\begin{practiceans}
  \correct{B}
  Balance the mass numbers: $27 + 4 = 30 + A$, so $A = 1$. Balance the atomic numbers:
  $13 + 2 = 15 + Z$, so $Z = 0$. A particle with $A = 1$ and no charge is a neutron,
  $\pracNuc{1}{0}{n}$.
  \begin{itemize}
    \item \xmark\ A: a proton is $\pracNuc{1}{1}{H}$. The $Z$ total on the right would be 16,
          not 15.
    \item \xmark\ C: an electron is $\pracNuc{0}{-1}{e}$. The $A$ total on the right would be
          30, not 31.
    \item \xmark\ D: a gamma ray has $A = 0$ and $Z = 0$. The $A$ total would be 30, not 31.
  \end{itemize}
\end{practiceans}

% ---- Q3 D+ meson
\Needspace{9\baselineskip}
\begin{practiceans}
  \correct{A}
  The charm quark is an up-type quark: $+\tfrac23 e$. The down quark has $-\tfrac13 e$, so
  the anti-down quark has the opposite sign: $+\tfrac13 e$. Total:
  $+\tfrac23 + \tfrac13 = \tfrac33 = +1e$. A quark and an antiquark make a meson. Two
  half-number spins add to a whole number, so a meson is a boson.
  \begin{itemize}
    \item \xmark\ B: fermions have half-number spin. Baryons (three quarks) are fermions;
          mesons are bosons.
    \item \xmark\ C: this forgets to flip the sign of the antiquark:
          $+\tfrac23 - \tfrac13 = +\tfrac13$.
    \item \xmark\ D: this makes both mistakes.
  \end{itemize}
\end{practiceans}

% ---- Q4 velocity selector with protons
\Needspace{9\baselineskip}
\begin{practiceans}
  \correct{D}
  The electrons go straight, so the electric force and the magnetic force are equal and
  opposite: $qE = qvB$, so $v = \dfrac{E}{B}$. The charge cancels, and the mass does not
  appear. A proton has the opposite charge, so \emph{both} forces point the other way. But
  their sizes, $eE$ and $evB$, are the same as before. They still cancel, so the protons go
  straight.
  \begin{itemize}
    \item \xmark\ A: the magnetic force on the protons is reversed too, so the two forces
          still balance.
    \item \xmark\ B: the net force is zero, so there is no bending. Mass does not appear in
          $qE = qvB$.
    \item \xmark\ C: the magnetic force $evB$ has the same size for protons as for electrons
          (same size of charge, same speed), and the electric force still balances it.
  \end{itemize}
\end{practiceans}

% ---- Q5 Chadwick
\Needspace{9\baselineskip}
\begin{practiceans}
  \correct{C}
  The fastest protons knocked out of the wax had about \qty{5.7}{MeV}. For its energy, a
  photon carries very little momentum. To give a proton \qty{5.7}{MeV}, a gamma photon would
  need about \qty{50}{MeV}. The alpha--beryllium reaction could supply at most about
  \qty{14}{MeV}. So gamma rays would break conservation of energy. A neutral particle with
  about a proton's mass fits the data.
  \begin{itemize}
    \item \xmark\ A: this is true, but gamma rays have no charge either. So it does not rule
          them out. (This is why Bothe and Becker thought it \emph{was} gamma rays.)
    \item \xmark\ B: false. The radiation was very penetrating: it passed through thick lead.
    \item \xmark\ D: neutrons leave no track, so their speed could not be measured directly.
          Chadwick worked out the mass from the recoil of the nuclei it hit.
  \end{itemize}
\end{practiceans}

% ---- Q6 longest wavelength
\Needspace{22\baselineskip}
\begin{practiceans}
  \correct{D}
  The longest wavelength comes from the smallest energy gap, because $E = \dfrac{hc}{\lambda}$.
  Higher levels are closer together. In the Rydberg equation, the smallest bracket
  $\left(\dfrac{1}{n_\mathrm{f}^2} - \dfrac{1}{n_\mathrm{i}^2}\right)$ gives the smallest
  $\dfrac{1}{\lambda}$, so the longest $\lambda$:
  \begin{center}\small
  \begin{tabular}{@{}lcccc@{}}
    \toprule
    Jump & $2 \to 1$ & $3 \to 2$ & $4 \to 2$ & $4 \to 3$\\
    \midrule
    Bracket & 0.7500 & 0.1389 & 0.1875 & 0.0486\\
    Energy gap (\si{eV}) & 10.2 & 1.89 & 2.55 & 0.661\\
    $\lambda$ (\si{nm}) & 121.5 (UV) & 656.3 (red) & 486.2 (cyan) & 1875 (infrared)\\
    \bottomrule
  \end{tabular}
  \end{center}
  \begin{itemize}
    \item \xmark\ A: the biggest gap gives the \emph{shortest} wavelength.
    \item \xmark\ B: \qty{656}{nm} is the longest \emph{visible} hydrogen line, but $4 \to 3$
          is longer (infrared).
    \item \xmark\ C: a bigger gap than $3 \to 2$, so a shorter wavelength.
  \end{itemize}
\end{practiceans}

% ---- Q7 binding energy curve
\Needspace{9\baselineskip}
\begin{practiceans}
  \correct{B}
  The binding energy per nucleon curve peaks near iron-56. Splitting a heavy nucleus and
  joining light nuclei both move the nucleons towards the peak. The products are more
  tightly bound, so their total mass is smaller. The lost mass comes out as energy:
  $E = \Delta m\,c^2$.
  \begin{itemize}
    \item \xmark\ A: binding energy is the energy needed to pull a nucleus apart. It is not
          stored energy that comes out. Here the products have a \emph{larger} value.
    \item \xmark\ C: the number of nucleons ($A$) is conserved. Instead, each nucleon ends up
          with a little less mass.
    \item \xmark\ D: in both reactions the products have \emph{less} mass. More mass would
          need energy to be supplied.
  \end{itemize}
\end{practiceans}

% ---- Q8 cyclotron
\Needspace{9\baselineskip}
\begin{practiceans}
  \correct{A}
  Each time an ion crosses the gap, the electric field does work on it ($W = qV$), so it
  speeds up. Inside the dees the magnetic force is at \ang{90} to the velocity. It bends the
  path into a half circle but does no work.
  \begin{itemize}
    \item \xmark\ B: a force at \ang{90} to the motion does no work, so the magnetic field
          cannot change the speed.
    \item \xmark\ C: $r = \dfrac{mv}{qB}$ grows as $v$ grows, so the path spirals outwards.
    \item \xmark\ D: the time for a half turn, $\dfrac{T}{2} = \dfrac{\pi m}{qB}$, does not
          depend on the speed (well below $c$). So one fixed frequency works.
  \end{itemize}
\end{practiceans}

% ---- Q9 magnitudes and radius
\Needspace{9\baselineskip}
\begin{practiceans}
  \correct{C}
  The difference is $7.5 - (-2.5) = 10$ magnitudes. That is two steps of 5, so the factor is
  $100 \times 100 = 10^4$. X has the smaller magnitude, so X is brighter. The stars have the
  same temperature, so each square metre of their surfaces gives out the same power. X must
  have $10^4$ times more surface area, so X is much bigger. (Area depends on radius$^2$, so its
  radius is $\sqrt{10^4} = 100$ times bigger.)
  \begin{itemize}
    \item \xmark\ A: a \emph{larger} magnitude means a \emph{dimmer} star.
    \item \xmark\ B: the magnitude scale is not linear. Ten magnitudes is a factor of $10^4$,
          not 10.
    \item \xmark\ D: if two stars had the same temperature and the same radius, they would
          have the same luminosity.
  \end{itemize}
\end{practiceans}

% ===================================================================== Section II
% ---- Q10 Millikan
\begin{practiceans}
  \given{$m = \qty{4.2e-15}{kg}$, \ $d = \qty{15.0}{mm} = \qty{0.0150}{m}$, \
    $V = \qty{964}{V}$, \ $g = \qty{9.8}{m.s^{-2}}$, \ $e = \qty{1.602e-19}{C}$}
  \find{(a) the size of the charge $q$ \quad (b) the number of extra electrons $n$}
  \formula{$q = \dfrac{mgd}{V}$ (not on the formula sheet: the drop is still, so
    $qE = mg$, with $E = \dfrac{V}{d}$) \quad and \quad $n = \dfrac{q}{e}$}
  \substitute{\eqlines{mgd &= \qty{4.2e-15}{kg} \times \qty{9.8}{m.s^{-2}} \times
      \qty{0.0150}{m}\\
    &= \qty{6.174e-16}{J} \quad\text{(top of the fraction)}\\
    q &= \frac{\qty{6.174e-16}{J}}{\qty{964}{V}} = \qty{6.405e-19}{C}
      \quad\text{(divide by $V$)}\\
    n &= \frac{\qty{6.405e-19}{C}}{\qty{1.602e-19}{C}} = 4.00 \quad\text{(count the electrons)}}}
  \answer{(a) \final{$q = \qty{6.4e-19}{C}$} \quad (b) \final{4 extra electrons}. The top
    plate is positive, so the field points down. The electric force must point up to hold
    the drop, so the drop is negative: it has gained electrons. $n$ is a whole number, as it
    must be.}
  \checkstep{units: $\unit{J} \div \unit{V} = \unit{J} \div \unit{J.C^{-1}} = \unit{C}$
    \cmark. A few electrons is typical for an oil drop \cmark.}
  \begin{markingcriteria}
    \criterion{(a) Uses $qE = mg$ with $E = \dfrac{V}{d}$, and $d$ in metres}{1}
    \criterion{(a) Correct charge, \qty{6.4e-19}{C}}{1}
    \criterion{(b) 4 electrons, justified (whole number of $e$; the drop is negative)}{1}
  \end{markingcriteria}
\end{practiceans}

% ---- Q11 oxygen-16 binding energy
\begin{practiceans}
  \given{$Z = 8$, \ $N = 16 - 8 = 8$, \ $m_\mathrm{nucleus} = \qty{15.990523}{\amu}$, \
    $m_\mathrm{p} = \qty{1.007276}{\amu}$, \ $m_\mathrm{n} = \qty{1.008665}{\amu}$}
  \find{(a) $\Delta m$ \quad (b) $E_\mathrm{B}$ \quad (c) $\dfrac{E_\mathrm{B}}{A}$}
  \formula{$\Delta m = (Zm_\mathrm{p} + Nm_\mathrm{n}) - m_\mathrm{nucleus}$ \ and \
    $E_\mathrm{B}\ (\si{MeV}) = \Delta m\ (\si{\amu}) \times 931.5$ (both not on the formula
    sheet; the second uses $E = mc^2$ with $\qty{1}{\amu} = \qty{931.5}{MeV}/c^2$)}
  \substitute{\eqlines{\text{parts} &= 8 \times \qty{1.007276}{\amu} + 8 \times
      \qty{1.008665}{\amu}\\
    &= \qty{8.058208}{\amu} + \qty{8.069320}{\amu} = \qty{16.127528}{\amu}\\
    \Delta m &= \qty{16.127528}{\amu} - \qty{15.990523}{\amu} = \qty{0.137005}{\amu}\\
    E_\mathrm{B} &= 0.137005 \times \qty{931.5}{MeV} = \qty{127.62}{MeV}\\
    \frac{E_\mathrm{B}}{A} &= \frac{\qty{127.62}{MeV}}{16} = \qty{7.976}{MeV}
      \quad\text{(16 nucleons)}}}
  \answer{(a) \final{$\Delta m = \qty{0.137005}{\amu}$} \quad
    (b) \final{$E_\mathrm{B} = \qty{127.6}{MeV}$}\par
    (c) \final{\qty{7.98}{MeV} per nucleon}\par
    (931.5 has 4 significant figures, so give $E_\mathrm{B}$ to 4 significant figures. In
    (c), divide the unrounded value by 16.)}
  \checkstep{the nucleus is lighter than its parts \cmark. About \qty{8}{MeV} per nucleon is
    near the top of the binding energy curve (iron-56 has \qty{8.8}{MeV}) \cmark.}
  \begin{markingcriteria}
    \criterion{(a) Correct mass defect, \qty{0.137005}{\amu}}{1}
    \criterion{(b) Correct binding energy, \qty{127.6}{MeV} (accept \qty{128}{MeV})}{1}
    \criterion{(c) Correct binding energy per nucleon, \qty{7.98}{MeV}}{1}
  \end{markingcriteria}
\end{practiceans}

% ---- Q12 Wien and spectra
\begin{practiceans}
  \given{$\lambda_\mathrm{max} = \qty{460}{nm} = \qty{4.60e-7}{m}$, \
    $b = \qty{2.898e-3}{m.K}$}
  \find{(a) the surface temperature $T$ \quad (b) the class and colour \quad (c) an
    explanation}
  \formula{$\lambda_\mathrm{max} = \dfrac{b}{T}$, so $T = \dfrac{b}{\lambda_\mathrm{max}}$}
  \substitute{$T = \dfrac{\qty{2.898e-3}{m.K}}{\qty{4.60e-7}{m}} = \qty{6300}{K}$\par
    The classes run O B A F G K M (hottest to coolest). Class F is \num{6000} to
    \qty{7500}{K}.}
  \answer{(a) \final{$T = \qty{6.30e3}{K}$} \quad (b) \final{class F}, yellow-white}
  \checkstep{a little hotter than the Sun (\qty{5800}{K}, peak near \qty{500}{nm}), so its
    peak is at a slightly shorter wavelength \cmark.}
  \begin{modelanswer}
  (c) The hot, dense inside of the star gives out a continuous (black body) spectrum. The
  cooler gas in the star's outer layers absorbs some photons. A photon is absorbed only if
  its energy exactly equals the gap between two energy levels ($hf = E_\mathrm{f} -
  E_\mathrm{i}$). The atoms give this energy out again in random directions, so less light
  of those wavelengths reaches us. Those wavelengths look dark. A black body spectrum is
  continuous: every wavelength is present, on a smooth curve whose peak depends only on the
  temperature. The star's spectrum is that curve with dark lines at wavelengths that identify
  the elements in its outer layers.
  \end{modelanswer}
  \begin{markingcriteria}
    \criterion{(a) Correct temperature, \qty{6.30e3}{K}}{1}
    \criterion{(b) Class F, yellow-white}{1}
    \criterion{(c) Explains absorption by cooler gas of photons whose energy matches a gap
      between energy levels, AND contrasts this with the continuous black body spectrum}{2}
    \criterion{(c) Gives one of these two points}{1}
  \end{markingcriteria}
\end{practiceans}

% ---- Q13 de Broglie: accelerating voltage
\begin{practiceans}
  \given{$\lambda = \qty{1.10e-10}{m}$, \ $m = \qty{9.109e-31}{kg}$, \
    $h = \qty{6.626e-34}{J.s}$, \ $q = \qty{1.602e-19}{C}$; the electrons start from rest}
  \find{(a) the speed $v$ \quad (b) the accelerating voltage $V$}
  \formula{$\lambda = \dfrac{h}{mv}$, so $v = \dfrac{h}{m\lambda}$; \ then
    $qV = \tfrac12 mv^2$ ($W = qV$ and $K = \tfrac12 mv^2$, both on the sheet), so
    $V = \dfrac{mv^2}{2q}$}
  \substitute{\eqlines{v &= \frac{\qty{6.626e-34}{J.s}}{\qty{9.109e-31}{kg} \times
      \qty{1.10e-10}{m}} = \qty{6.613e6}{m.s^{-1}}\\
    K &= \tfrac12 \times \qty{9.109e-31}{kg} \times (\qty{6.613e6}{m.s^{-1}})^2
      = \qty{1.992e-17}{J}\\
    V &= \frac{K}{q} = \frac{\qty{1.992e-17}{J}}{\qty{1.602e-19}{C}} = \qty{124.3}{V}}}
  \answer{(a) \final{$v = \qty{6.61e6}{m.s^{-1}}$} \quad (b) \final{$V = \qty{124}{V}$}}
  \checkstep{$\dfrac{v}{c} = \dfrac{\num{6.61e6}}{\num{3.00e8}} = 0.022$, well below $0.1c$,
    so $K = \tfrac12 mv^2$ is fine. Davisson and Germer used a similar voltage (\qty{54}{V})
    \cmark.}
  \begin{markingcriteria}
    \criterion{(a) Correct speed, \qty{6.61e6}{m.s^{-1}}}{1}
    \criterion{(b) Uses $qV = \tfrac12 mv^2$}{1}
    \criterion{(b) Correct voltage, \qty{124}{V}}{1}
  \end{markingcriteria}
\end{practiceans}

% ---- Q14 redshift and Hubble's law
\begin{practiceans}
  \given{$\lambda_0 = \qty{486.1}{nm}$, \ $\lambda = \qty{497.8}{nm}$, \
    $H_0 = \qty{68.0}{km.s^{-1}.Mpc^{-1}}$, \ $c = \qty{3.00e8}{m.s^{-1}}$}
  \find{(a) the recession velocity $v$ \quad (b) the distance $d$ \quad (c) an explanation}
  \formula{$\dfrac{\Delta\lambda}{\lambda_0} \approx \dfrac{v}{c}$, so
    $v = \dfrac{c\,\Delta\lambda}{\lambda_0}$; \ $v = H_0 d$, so $d = \dfrac{v}{H_0}$ (both not
    on the formula sheet)}
  \substitute{\eqlines{\Delta\lambda &= \qty{497.8}{nm} - \qty{486.1}{nm} = \qty{11.7}{nm}\\
    \frac{\Delta\lambda}{\lambda_0} &= \frac{\qty{11.7}{nm}}{\qty{486.1}{nm}} = 0.02407
      \quad\text{(no unit)}\\
    v &= 0.02407 \times \qty{3.00e8}{m.s^{-1}} = \qty{7.221e6}{m.s^{-1}}\\
    &= \qty{7221}{km.s^{-1}} \quad\text{(divide by 1000 for km)}\\
    d &= \frac{\qty{7221}{km.s^{-1}}}{\qty{68.0}{km.s^{-1}.Mpc^{-1}}} = \qty{106.2}{Mpc}}\par
    \textbf{(c)} Space itself is expanding, everywhere. An observer in \emph{any} galaxy
    would see almost all other galaxies moving away, with speed proportional to distance.
    So no galaxy is special, and there is no centre.}
  \answer{(a) \final{$v = \qty{7.22e6}{m.s^{-1}}$} ($\qty{7.22e3}{km.s^{-1}}$) \quad
    (b) \final{$d = \qty{106}{Mpc}$}\par
    (c) every galaxy sees the same pattern, so the Milky Way is not at the centre.}
  \checkstep{$v$ is only \qty{2.4}{\percent} of $c$, so the low-speed redshift formula is
    fine \cmark. Galaxies about \qty{100}{Mpc} away move at several thousand
    \si{km.s^{-1}} \cmark.}
  \begin{markingcriteria}
    \criterion{(a) Finds $\dfrac{\Delta\lambda}{\lambda_0} = 0.0241$}{1}
    \criterion{(a) Correct velocity, \qty{7.22e6}{m.s^{-1}}}{1}
    \criterion{(b) Correct distance, \qty{106}{Mpc}}{1}
    \criterion{(c) Space expands everywhere, so observers in every galaxy see the same
      recession (no centre)}{1}
  \end{markingcriteria}
\end{practiceans}

% ---- Q15 proton-proton chain step
\begin{practiceans}
  \given{$m(\mathrm{H\text{-}2}) = \qty{2.014102}{\amu}$, \
    $m(\mathrm{H\text{-}1}) = \qty{1.007825}{\amu}$, \
    $m(\mathrm{He\text{-}3}) = \qty{3.016029}{\amu}$. Electrons balance ($1 + 1 = 2$), so
    the atomic masses can be used as they are.}
  \find{(a) the energy released $E$, in MeV \quad (b) the photon's wavelength $\lambda$}
  \formula{$E\ (\si{MeV}) = (m_\mathrm{before} - m_\mathrm{after}) \times 931.5$ (from
    $E = mc^2$); \ $E = \dfrac{hc}{\lambda}$ (not on the sheet: $E = hf$ with $c = f\lambda$),
    so $\lambda = \dfrac{hc}{E}$}
  \substitute{\eqlines{m_\mathrm{before} &= \qty{2.014102}{\amu} + \qty{1.007825}{\amu}
      = \qty{3.021927}{\amu}\\
    \Delta m &= \qty{3.021927}{\amu} - \qty{3.016029}{\amu} = \qty{0.005898}{\amu}\\
    E &= 0.005898 \times \qty{931.5}{MeV} = \qty{5.494}{MeV}\\
    E &= 5.494 \times \qty{1.602e-13}{J} = \qty{8.801e-13}{J} \quad\text{(MeV to J)}\\
    \lambda &= \frac{\qty{6.626e-34}{J.s} \times \qty{3.00e8}{m.s^{-1}}}{\qty{8.801e-13}{J}}
      = \qty{2.26e-13}{m}}}
  \answer{(a) \final{$E = \qty{5.49}{MeV}$} \quad (b) \final{$\lambda = \qty{2.26e-13}{m}$}}
  \checkstep{a gamma ray has a wavelength much shorter than \qty{e-11}{m} \cmark. The energy
    is positive because the product has less mass \cmark.}
  \begin{markingcriteria}
    \criterion{(a) Correct mass defect, \qty{0.005898}{\amu}}{1}
    \criterion{(a) Correct energy, \qty{5.49}{MeV}}{1}
    \criterion{(b) Changes MeV to J and finds $\lambda = \qty{2.26e-13}{m}$}{1}
  \end{markingcriteria}
\end{practiceans}

% ---- Q16 Thomson
\begin{practiceans}
  \given{plates $d = \qty{2.0}{cm} = \qty{0.020}{m}$, \ $V = \qty{400}{V}$, \
    $B = \qty{8.0e-4}{T}$, \ $r = \qty{17.8}{cm} = \qty{0.178}{m}$, \
    $e = \qty{1.602e-19}{C}$}
  \find{(a) the speed $v$ \quad (b) $\dfrac{q}{m}$ \quad (c) the mass $m_\mathrm{e}$}
  \formula{$E = \dfrac{V}{d}$ (on the sheet); \ $v = \dfrac{E}{B}$ (not on the sheet:
    $qE = qvB$); \ $\dfrac{q}{m} = \dfrac{v}{Br}$ (not on the sheet: $qvB = \dfrac{mv^2}{r}$);
    \ $m_\mathrm{e} = e \div \dfrac{q}{m}$}
  \substitute{\eqlines{E &= \frac{\qty{400}{V}}{\qty{0.020}{m}} = \qty{2.0e4}{V.m^{-1}}\\
    v &= \frac{\qty{2.0e4}{V.m^{-1}}}{\qty{8.0e-4}{T}} = \qty{2.5e7}{m.s^{-1}}\\
    Br &= \qty{8.0e-4}{T} \times \qty{0.178}{m} = \num{1.424e-4} \quad\text{(bottom of the
      fraction)}\\
    \frac{q}{m} &= \frac{\qty{2.5e7}{m.s^{-1}}}{\num{1.424e-4}} = \qty{1.756e11}{C.kg^{-1}}\\
    m_\mathrm{e} &= \frac{\qty{1.602e-19}{C}}{\qty{1.756e11}{C.kg^{-1}}} = \qty{9.12e-31}{kg}}}
  \answer{(a) \final{$v = \qty{2.5e7}{m.s^{-1}}$} \quad
    (b) \final{$\dfrac{q}{m} = \qty{1.8e11}{C.kg^{-1}}$}\par
    (c) \final{$m_\mathrm{e} = \qty{9.1e-31}{kg}$} (2 significant figures, like $B$)\par
    In (c), divide by the unrounded $\dfrac{q}{m}$ (\num{1.756e11}). The rounded value
    \num{1.8e11} gives \qty{8.9e-31}{kg}, which is too small.}
  \checkstep{the data-sheet value is \qty{9.109e-31}{kg} \cmark. The speed is about
    \qty{8}{\percent} of $c$, so relativity can be ignored \cmark.}
  \begin{markingcriteria}
    \criterion{(a) Correct field, \qty{2.0e4}{V.m^{-1}}}{1}
    \criterion{(a) Correct speed, \qty{2.5e7}{m.s^{-1}}}{1}
    \criterion{(b) Correct charge-to-mass ratio, \qty{1.8e11}{C.kg^{-1}}}{1}
    \criterion{(c) Correct mass, \qty{9.1e-31}{kg}}{1}
  \end{markingcriteria}
\end{practiceans}

% ---- Q17 technetium-99m
\begin{practiceans}
  \given{$t_{\frac{1}{2}} = \qty{6.01}{h} = 6.01 \times \qty{3600}{s} = \qty{21636}{s}$, \
    $A_0 = \qty{6.0e8}{Bq}$, \ $t = \qty{24.0}{h} = \qty{86400}{s}$}
  \find{(a) $\lambda$ \quad (b) $N_0$ \quad (c) $A_\mathrm{t}$ after \qty{24.0}{h} \quad
    (d) the time for $A_\mathrm{t} = 0.010A_0$}
  \formula{$\lambda = \dfrac{\ln 2}{t_{\frac{1}{2}}}$ (on the sheet); \ $A = \lambda N$ and
    $A_\mathrm{t} = A_0 e^{-\lambda t}$ (not on the sheet; the same form as
    $N_\mathrm{t} = N_0 e^{-\lambda t}$)}
  \substitute{\eqlines{\lambda &= \frac{0.6931}{\qty{21636}{s}} = \qty{3.204e-5}{s^{-1}}\\
    N_0 &= \frac{A_0}{\lambda} = \frac{\qty{6.0e8}{s^{-1}}}{\qty{3.204e-5}{s^{-1}}}
      = \num{1.873e13}\\
    \lambda t &= \qty{3.204e-5}{s^{-1}} \times \qty{86400}{s} = 2.768 \quad\text{(no unit)}\\
    A_\mathrm{t} &= \qty{6.0e8}{Bq} \times e^{-2.768} = \qty{6.0e8}{Bq} \times 0.06279
      = \qty{3.77e7}{Bq}}\par
    \textbf{(d)} Divide $A_\mathrm{t} = A_0 e^{-\lambda t}$ by $A_0$, then take ln of both
    sides:\par
    \eqlines{0.010 &= e^{-\lambda t} \ \Rightarrow\ \ln 0.010 = -\lambda t
      \ \Rightarrow\ -4.605 = -\lambda t\\
    t &= \frac{4.605}{\qty{3.204e-5}{s^{-1}}} = \qty{1.437e5}{s}
      = \frac{\qty{1.437e5}{s}}{\qty{3600}{s.h^{-1}}} = \qty{39.9}{h}}}
  \answer{(a) \final{$\lambda = \qty{3.20e-5}{s^{-1}}$} \quad
    (b) \final{$N_0 = \num{1.9e13}$ nuclei}\par
    (c) \final{$A = \qty{3.8e7}{Bq}$} \quad (d) \final{$t = \qty{39.9}{h}$} (about
    \qty{40}{h})}
  \checkstep{\qty{24.0}{h} is almost 4 half-lives, and $\left(\tfrac12\right)^4 = 0.0625$
    \cmark. Falling to \qty{1}{\percent} takes $\dfrac{\ln 100}{\ln 2} = 6.6$ half-lives,
    and $6.6 \times \qty{6.01}{h} = \qty{40}{h}$ \cmark.}
  \begin{markingcriteria}
    \criterion{(a) Correct decay constant in \si{s^{-1}}}{1}
    \criterion{(b) Correct number of nuclei, using $A = \lambda N$}{1}
    \criterion{(c) Correct activity, \qty{3.8e7}{Bq}}{1}
    \criterion{(d) Correct time, \qty{39.9}{h} (accept \qty{40}{h})}{1}
  \end{markingcriteria}
\end{practiceans}

% ---- Q18 fission
\begin{practiceans}
  \given{the fission equation; binding energy per nucleon: U-235 \qty{7.59}{MeV}, Ba-141
    \qty{8.33}{MeV}, Kr-92 \qty{8.51}{MeV}; a free neutron has $E_\mathrm{B} = 0$}
  \find{(a) $x$ \quad (b) the energy released $E$ \quad (c) the job of the control rods}
  \formula{$\sum A$ and $\sum Z$ are the same on both sides; \
    $E_\mathrm{B} = (\text{per nucleon value}) \times A$; \
    $E = \sum E_\mathrm{B,after} - \sum E_\mathrm{B,before}$ (not on the sheet)}
  \substitute{\eqlines{A\colon\ 235 + 1 &= 141 + 92 + x \ \Rightarrow\ x = 236 - 233 = 3\\
    Z\colon\ 92 + 0 &= 56 + 36 + 0 = 92 \ \text{\cmark}\\
    E_\mathrm{B,before} &= 235 \times \qty{7.59}{MeV} = \qty{1783.65}{MeV}\\
    E_\mathrm{B,after} &= 141 \times \qty{8.33}{MeV} + 92 \times \qty{8.51}{MeV}\\
    &= \qty{1174.53}{MeV} + \qty{782.92}{MeV} = \qty{1957.45}{MeV}\\
    E &= \qty{1957.45}{MeV} - \qty{1783.65}{MeV} = \qty{173.8}{MeV}}\par
    \textbf{(c)} Control rods (boron or cadmium) absorb neutrons. They are moved in or out
    so that, on average, exactly one neutron from each fission causes another fission. Then
    the rate of fission stays steady.}
  \answer{(a) \final{$x = 3$} \quad (b) \final{$E \approx \qty{174}{MeV}$} \quad (c) they
    absorb neutrons, so that on average one neutron per fission causes a new fission.}
  \checkstep{close to the usual ``about \qty{200}{MeV}'' per fission \cmark. Using atomic
    masses instead gives \qty{173}{MeV}, which agrees.}
  \begin{markingcriteria}
    \criterion{(a) $x = 3$}{1}
    \criterion{(b) Multiplies each per nucleon value by its mass number}{1}
    \criterion{(b) Correct energy, about \qty{174}{MeV}}{1}
    \criterion{(c) Control rods absorb neutrons, so on average one neutron per fission causes
      another fission}{1}
  \end{markingcriteria}
\end{practiceans}

% ---- Q19 deuteron cyclotron
\begin{practiceans}
  \given{$q = \qty{1.602e-19}{C}$, \ $m = \qty{3.344e-27}{kg}$, \ $B = \qty{1.2}{T}$, \
    $r = \qty{0.45}{m}$, \ gap voltage $V = \qty{40}{kV} = \qty{4.0e4}{V}$}
  \find{(a) the AC frequency $f$ \quad (b) the final kinetic energy $K$ \quad (c) the number
    of turns}
  \formula{$f = \dfrac{qB}{2\pi m}$ and $r = \dfrac{mv}{qB}$, so $v = \dfrac{qBr}{m}$ (both
    not on the sheet; from $qvB = \dfrac{mv^2}{r}$ and $v = \dfrac{2\pi r}{T}$); \
    $K = \tfrac12 mv^2$; \ each turn crosses the gap twice, gaining $2qV$}
  \substitute{\eqlines{f &= \frac{\qty{1.602e-19}{C} \times \qty{1.2}{T}}{2\pi \times
      \qty{3.344e-27}{kg}} = \qty{9.149e6}{Hz}\\
    v &= \frac{\qty{1.602e-19}{C} \times \qty{1.2}{T} \times \qty{0.45}{m}}{\qty{3.344e-27}{kg}}
      = \qty{2.587e7}{m.s^{-1}}\\
    K &= \tfrac12 \times \qty{3.344e-27}{kg} \times (\qty{2.587e7}{m.s^{-1}})^2\\
    &= \qty{1.119e-12}{J}\\
    K &= \frac{\qty{1.119e-12}{J}}{\qty{1.602e-13}{J.MeV^{-1}}} = \qty{6.985}{MeV}
      \quad\text{(J to MeV)}\\
    \text{energy per turn} &= 2 \times \qty{40}{keV} = \qty{80}{keV} = \qty{0.080}{MeV}\\
    \text{turns} &= \frac{\qty{6.985}{MeV}}{\qty{0.080}{MeV}} = 87.3}}
  \answer{(a) \final{$f = \qty{9.1e6}{Hz}$} \quad (b) \final{$K = \qty{7.0}{MeV}$} \quad
    (c) \final{about 87 turns}}
  \checkstep{$\dfrac{v}{c} = \dfrac{\num{2.587e7}}{\num{3.00e8}} = 0.086$, below $0.1c$, so
    $\tfrac12 mv^2$ is fine \cmark. A bigger gap voltage would mean fewer turns, but the same
    final energy (it is set by $B$ and $r$) \cmark.}
  \begin{markingcriteria}
    \criterion{(a) Correct frequency, \qty{9.1e6}{Hz}}{1}
    \criterion{(b) Finds the exit speed from $r = \dfrac{mv}{qB}$}{1}
    \criterion{(b) Correct kinetic energy, \qty{7.0}{MeV}}{1}
    \criterion{(c) About 87 turns (two gap crossings per turn)}{1}
  \end{markingcriteria}
\end{practiceans}

% ===================================================================== Section III
% ---- Q20 Geiger-Marsden
% layout QA (book build): keep "Question 20" with the start of its model answer
\Needspace{9\baselineskip}
\begin{practiceans}
  \begin{modelanswer}
  \textbf{(a)} \emph{Results.} Most alpha particles passed straight through the gold foil,
  or were turned by only a small angle. A very small fraction (about 1 in 8000) were turned
  by more than \ang{90}. A few bounced almost straight back.

  \emph{Rutherford's explanation.}
  \begin{itemize}
    \item Most alphas passed straight through, so the atom is mostly empty space.
    \item A large deflection needs a very large repulsive force. Positive charge spread
          through the whole atom (Thomson's model) is far too weak to give it. So all the
          positive charge, and nearly all the mass, must be packed into a tiny centre: the
          nucleus.
    \item Large deflections were very rare, so the nucleus must be very small: about
          \num{10000} to \num{100000} times smaller across than the atom.
    \item The electrons move around the nucleus, far away from it, so the atom is neutral.
  \end{itemize}

  \textbf{(b)} An electron moving around the nucleus in a circle is accelerating (its
  direction keeps changing). Classical electromagnetism (Maxwell) says that an accelerating
  charge gives out electromagnetic waves. These carry energy away. So the electron should
  lose energy and spiral into the nucleus in about \qty{e-11}{s}. Real atoms are stable,
  so the model cannot be complete.
  \end{modelanswer}
  \begin{markingcriteria}
    \criterion{(a) Describes the results, including the rare large-angle deflections, AND
      links each result to a feature of the nuclear model (empty space; tiny, dense, positive
      nucleus)}{3}
    \criterion{(a) Describes the results and links one of them to the nuclear model}{2}
    \criterion{(a) States one correct result or one feature of the nuclear model}{1}
    \criterion{(b) The orbiting electron accelerates, so it radiates energy, so it spirals
      into the nucleus}{2}
    \criterion{(b) States that the electron radiates energy, OR that the atom should
      collapse}{1}
  \end{markingcriteria}
\end{practiceans}

% ---- Q21 Balmer investigation
\begin{practiceans}
  \given{500 lines per mm, so $d = \dfrac{\qty{1}{mm}}{500} = \dfrac{\qty{1.00e-3}{m}}{500}
    = \qty{2.00e-6}{m}$; \ $m = 1$; \ $\theta = \ang{14.1}$; \ $R = \qty{1.097e7}{m^{-1}}$}
  \find{(a) $\lambda$ \quad (b) the Rydberg wavelength for $n = 4 \to n = 2$ \quad
    (c) an evaluation}
  \formula{$d\sin\theta = m\lambda$ (on the sheet), so $\lambda = d\sin\theta$ for the first
    order; \ $\dfrac{1}{\lambda} = R\left(\dfrac{1}{n_\mathrm{f}^2} -
    \dfrac{1}{n_\mathrm{i}^2}\right)$ (on the sheet)}
  \substitute{\eqlines{\lambda &= \qty{2.00e-6}{m} \times \sin\ang{14.1}
      = \qty{2.00e-6}{m} \times 0.2436\\
    &= \qty{4.87e-7}{m} = \qty{487}{nm}\\
    \frac{1}{\lambda} &= \qty{1.097e7}{m^{-1}} \times \left(\frac{1}{2^2} - \frac{1}{4^2}\right)\\
    &= \qty{1.097e7}{m^{-1}} \times 0.1875 = \qty{2.057e6}{m^{-1}}\\
    \lambda &= \frac{1}{\qty{2.057e6}{m^{-1}}} = \qty{4.862e-7}{m} = \qty{486.2}{nm}}}
  \answer{(a) \final{$\lambda = \qty{487}{nm}$} \quad (b) the $4 \to 2$ jump gives
    \qty{486.2}{nm}. This is within \qty{0.2}{\percent} of the measured value. The other
    visible Balmer lines are far away (\qty{656}{nm}, \qty{434}{nm}, \qty{410}{nm}). So the
    cyan line is the $n = 4 \to n = 2$ jump.}
  \checkstep{cyan light is about 480 to \qty{500}{nm} \cmark.}
  \begin{modelanswer}
  (c) The method has low reliability. There is only one reading, so a random error (in
  setting the cross-wire on the line, or in reading the scale) cannot be seen or averaged
  out. Reading on one side only also lets a systematic error into $\theta$: for example, if
  the scale does not read exactly \ang{0} for the straight-through beam, or the grating is
  not at right angles to the beam. An error of only \ang{0.1} changes $\lambda$ by about
  \qty{3}{nm}. So the close agreement with the Rydberg value may be partly luck.

  \emph{Improvements:}
  \begin{enumerate}
    \item Repeat the reading at least three times and use the mean (better reliability).
    \item Measure the angle on both sides of the straight-through position, and halve the
          angle between them. This cancels a zero error and a slightly turned grating (better
          accuracy).
  \end{enumerate}
  \end{modelanswer}
  \begin{markingcriteria}
    \criterion{(a) Correct grating spacing, \qty{2.00e-6}{m}}{1}
    \criterion{(a) Correct wavelength, \qty{487}{nm}}{1}
    \criterion{(b) Rydberg value for $4 \to 2$ (\qty{486}{nm}) compared with the measured
      value}{1}
    \criterion{(c) Judges the reliability (low) with a reason, AND gives two suitable
      improvements}{3}
    \criterion{(c) Judges the reliability with a reason AND gives one suitable improvement,
      OR gives two suitable improvements}{2}
    \criterion{(c) Gives one suitable improvement, OR a judgement with a reason}{1}
  \end{markingcriteria}
\end{practiceans}

% ---- Q22 decay graph
\begin{practiceans}
  \textbf{(a)} Take ln of both sides of $N_\mathrm{t} = N_0 e^{-\lambda t}$:
  \begin{align*}
    \ln N_\mathrm{t} &= \ln\left(N_0 e^{-\lambda t}\right) && \text{take ln of both sides}\\
    \ln N_\mathrm{t} &= \ln N_0 + \ln\left(e^{-\lambda t}\right) && \text{$\ln(ab) = \ln a + \ln b$}\\
    \ln N_\mathrm{t} &= -\lambda t + \ln N_0 && \text{$\ln(e^{x}) = x$}
  \end{align*}
  This has the form $y = mx + c$, with $y = \ln N_\mathrm{t}$ and $x = t$. So the graph is a
  straight line. Its gradient is $-\lambda$ and its $y$-intercept is $\ln N_0$.

  \textbf{(b)} Corrected count rate = measured count rate $-$ 30. Then press ln.
  \begin{center}\small\setlength\tabcolsep{4pt}
  \begin{tabular}{@{}L{36mm}*{9}{C{8.5mm}}@{}}
    \toprule
    Time $t$ (\si{s}) & 0 & 30 & 60 & 90 & 120 & 150 & 180 & 210 & 240\\
    \midrule
    Corrected count rate & 971 & 767 & 547 & 407 & 298 & 220 & 165 & 120 & 87\\
    \midrule
    $\ln(\text{corrected})$ & 6.88 & 6.64 & 6.30 & 6.01 & 5.70 & 5.39 & 5.11 & 4.79 & 4.47\\
    \bottomrule
  \end{tabular}
  \end{center}

  \textbf{(c)} The points lie close to a straight line (see the graph below). Draw one
  straight line with about as many points above it as below it.
  \begin{figure}
    \centering
    \begin{tikzpicture}
      \begin{axis}[hsc, width=0.8\linewidth, height=0.52\linewidth,
          xmin=0, xmax=250, ymin=4.0, ymax=7.0,
          xlabel={Time $t$ (\si{s})}, ylabel={$\ln(\text{corrected count rate})$},
          xtick={0,50,...,250}, ytick={4.0,4.5,...,7.0},
          minor x tick num=4, minor y tick num=4,
          grid=both, minor grid style={hscrule!45, line width=0.2pt},
          yticklabel style={/pgf/number format/.cd, fixed, fixed zerofill, precision=1}]
        \addplot[hscseriesB, line width=1pt, domain=0:250, samples=2] {6.913 - 0.010125*x};
        \addplot[hscseriesA, only marks, mark=*, mark size=2pt] coordinates {
          (0,6.878) (30,6.642) (60,6.304) (90,6.009) (120,5.697) (150,5.394) (180,5.106)
          (210,4.787) (240,4.466)};
        \draw[hscmuted, line width=0.7pt, dashed] (axis cs:0,6.91) -- (axis cs:240,6.91)
          -- (axis cs:240,4.48);
        \node[font=\footnotesize, anchor=north, text=hscink] at (axis cs:165,6.89)
          {run $= \qty{240}{s}$};
        \node[font=\footnotesize, anchor=east, text=hscink] at (axis cs:237,5.9)
          {rise $= -2.43$};
        \node[font=\footnotesize, anchor=north east, text=hscseriesB] at (axis cs:120,5.6)
          {line of best fit};
      \end{axis}
    \end{tikzpicture}
    \caption*{Model graph for part (c). The dashed triangle gives the gradient in part (d).}
  \end{figure}

  \textbf{(d)}
  \given{two points far apart \emph{on the line}: $(\qty{0}{s},\ 6.91)$ and
    $(\qty{240}{s},\ 4.48)$}
  \find{the gradient, then $\lambda$, then $t_{\frac{1}{2}}$}
  \formula{gradient $= \dfrac{\text{rise}}{\text{run}} = -\lambda$; \
    $\lambda = \dfrac{\ln 2}{t_{\frac{1}{2}}}$ (on the sheet), so
    $t_{\frac{1}{2}} = \dfrac{\ln 2}{\lambda}$}
  \substitute{\eqlines{\text{gradient} &= \frac{4.48 - 6.91}{\qty{240}{s} - \qty{0}{s}}
      = \frac{-2.43}{\qty{240}{s}} = \qty{-0.01013}{s^{-1}}\\
    \lambda &= \qty{0.01013}{s^{-1}} \quad\text{(gradient $= -\lambda$)}\\
    t_{\frac{1}{2}} &= \frac{0.6931}{\qty{0.01013}{s^{-1}}} = \qty{68.4}{s}}}
  \answer{\final{$t_{\frac{1}{2}} \approx \qty{68}{s}$}. (Any value from about 65 to
    \qty{72}{s} is fine; it depends on your line.)}
  \checkstep{from the table, the count rate falls to a quarter of 971 (about 243) at about
    \qty{140}{s}. That is two half-lives, so $t_{\frac{1}{2}} \approx \qty{70}{s}$, close to
    our answer \cmark. The accepted value is about \qty{70}{s}. The points scatter a little
    about the line because decay is random.}
  \begin{markingcriteria}
    \criterion{(a) Takes ln of both sides and shows $\ln N_\mathrm{t} = \ln N_0 - \lambda t$,
      a straight line with gradient $-\lambda$}{1}
    \criterion{(b) All corrected count rates and ln values correct}{1}
    \criterion{(c) All 9 points plotted correctly, AND a single straight line of best
      fit}{2}
    \criterion{(c) At least 7 points plotted correctly with a suitable line, OR all points
      plotted correctly but the line is missing or poor}{1}
    \criterion{(d) Gradient found from two points on the line}{1}
    \criterion{(d) Half-life from $t_{\frac{1}{2}} = \dfrac{\ln 2}{\lambda}$, about
      \qty{68}{s}}{1}
  \end{markingcriteria}
\end{practiceans}

% ---- Q23 matter waves: Bohr to Schrodinger
\begin{practiceans}
  \begin{modelanswer}
  \textbf{Bohr's limitation.} Bohr's model (1913) said that electrons move only in certain
  orbits (stationary states), where they do not radiate. Their angular momentum is
  quantised: $mvr = \dfrac{nh}{2\pi}$. But these were \emph{postulates}. Bohr gave no reason
  why only these orbits are allowed, or why an electron in them does not radiate.

  \textbf{De Broglie's matter waves (1924).} De Broglie proposed that every moving particle
  has a wavelength, $\lambda = \dfrac{h}{mv}$. Then an electron in an orbit is a wave
  wrapped around the nucleus. The wave joins itself smoothly only if a whole number of
  wavelengths fits around the orbit: $2\pi r = n\lambda$. This makes a standing wave. In any
  other orbit the wave is out of step with itself and cancels, so that orbit cannot exist.
  Put $\lambda = \dfrac{h}{mv}$ into $2\pi r = n\lambda$ and you get $mvr = \dfrac{nh}{2\pi}$:
  exactly Bohr's rule. A standing wave does not travel, so it does not carry energy away. So
  matter waves \emph{explain} what Bohr had only assumed.

  \textbf{Evidence.} In 1927 Davisson and Germer, and G.\,P. Thomson, diffracted electrons
  with crystals. The wavelengths from the patterns matched $\dfrac{h}{mv}$. So matter waves
  are real.

  \textbf{Schrödinger's model (1926).} Schrödinger treated the electron fully as a wave. He
  wrote a wave equation for it in three dimensions. Its solutions (wave functions) describe
  orbitals: regions where the electron is \emph{likely} to be found. There are no fixed
  orbits. The energy levels and three quantum numbers ($n$, $l$, $m_l$) come out of the maths;
  none is assumed. With approximations, the model also works for atoms with many electrons.
  It also predicts how bright the spectral lines are.

  \textbf{Analysis.} De Broglie's idea is the link between the two models. It turned Bohr's
  assumed rule into a result of wave behaviour. It also gave Schrödinger the starting point
  for today's quantum model of the atom.
  \end{modelanswer}
  \begin{markingcriteria}
    \criterion{Identifies a limitation of Bohr's model; explains how matter waves
      ($\lambda = \dfrac{h}{mv}$, $2\pi r = n\lambda$) explain the stationary states; refers
      to evidence (electron diffraction); explains how Schrödinger's wave model developed
      from this (wave functions, probability, quantisation from the maths); clear
      cause-and-effect links throughout}{5}
    \criterion{Covers most of these points, with one part weak or links missing}{4}
    \criterion{Explains de Broglie's idea and links it to Bohr's model OR to Schrödinger's
      model}{3}
    \criterion{Describes some relevant features of the models or of matter waves}{2}
    \criterion{States one relevant point}{1}
  \end{markingcriteria}
\end{practiceans}

% ---- Q24 accelerators as microscopes
\begin{practiceans}
  \begin{modelanswer}
  \textbf{How accelerators work.} Electric fields speed up charged particles: in the gaps
  of a linac, in the gap between the dees of a cyclotron, and in the RF cavities of a
  synchrotron. Magnetic fields bend and focus the beam ($r = \dfrac{p}{qB}$). In a
  synchrotron the magnetic field is raised as the momentum grows, so the particles stay on
  one ring. In a collider two beams meet head-on, and detectors track the particles that
  come out.

  \textbf{Why ``microscope'' fits.} To see a small detail, the probe's wavelength must be
  smaller than the detail. A faster particle has more momentum, so its de Broglie wavelength,
  $\lambda = \dfrac{h}{p}$, is shorter. Accelerators give particles enough momentum to make
  $\lambda$ smaller than a proton (about \qty{e-15}{m}). At SLAC (1967--1973), electrons of up
  to about \qty{20}{GeV} ($\lambda \approx \qty{6e-17}{m}$) were fired at protons. Far more
  electrons than expected were scattered through large angles. So the proton contains tiny,
  point-like charges: quarks. Protons are not fundamental. This is Rutherford's gold-foil
  idea, one level deeper.

  \textbf{Why the description is incomplete.} Accelerators also \emph{create} particles that
  were not inside the colliding particles. The collision energy becomes mass ($E = mc^2$).
  The Standard Model predicted particles, and accelerators then found them with the
  predicted properties: the $\Omega^-$ (1964), the W and Z bosons at about 80 and
  \qty{91}{GeV}/$c^2$ (CERN, 1983), the top quark (Fermilab, 1995) and the Higgs boson at
  \qty{125}{GeV}/$c^2$ (LHC, 2012). A microscope only looks at what is already there.

  \textbf{Judgement.} The description is partly valid. Accelerators act like microscopes when
  they probe structure, as in deep inelastic scattering. But it is incomplete: much of their
  role in testing the Standard Model is as ``particle factories'', which turn energy into
  new, heavy particles.
  \end{modelanswer}
  \begin{markingcriteria}
    \criterion{Explains how accelerators work (electric fields accelerate, magnetic fields
      steer); explains the microscope idea with the de Broglie wavelength, with evidence (for
      example deep inelastic scattering); explains particle creation with $E = mc^2$, with
      evidence (for example W and Z, Higgs); makes a clear, justified judgement}{6}
    \criterion{As above, with one part less detailed}{5}
    \criterion{Explains how accelerators work and gives two pieces of evidence with a
      judgement; OR explains both roles without specific evidence}{4}
    \criterion{Describes how an accelerator works and links one piece of evidence to the
      Standard Model}{3}
    \criterion{Describes some features of accelerators, or some evidence}{2}
    \criterion{States one relevant point}{1}
  \end{markingcriteria}
\end{practiceans}

% ---- Q25 HR diagram
% layout QA (book build): one extra line on this page and on the next lets the Q25
% marking criteria end on the same page as the model answer.
\enlargethispage{\baselineskip}
\begin{practiceans}
  \begin{modelanswer}
  \textbf{(a)} \emph{R is a white dwarf.} It is hot (about \qty{9000}{K}, on the left) but
  very dim (about $0.003\,L_\odot$, near the bottom). A hot star gives out a lot of power
  from each square metre, so a dim hot star must be very small.

  \emph{P is a red giant.} It is cool (about \qty{3600}{K}, on the right, so red) but very
  luminous (about $300\,L_\odot$). A cool star gives out little power from each square metre,
  so a luminous cool star must be very large. It lies above and to the right of the main
  sequence.

  \textbf{(b)}
  \begin{itemize}
    \item \emph{Mass:} on the main sequence, hotter and more luminous stars are more massive.
          Q (about \qty{20000}{K}, $3000\,L_\odot$) is much more massive than the Sun. S
          (about \qty{4200}{K}, $0.15\,L_\odot$) is less massive than the Sun.
    \item \emph{Time on the main sequence:} Q has more fuel, but its hotter core fuses it
          much faster. So Q stays on the main sequence for only tens of millions of years.
          S stays there for tens of billions of years, longer than the Sun.
    \item \emph{Fusion:} both fuse hydrogen into helium in their cores. Q's core is hotter
          than about 18 million K, so Q mainly uses the CNO cycle. Carbon, nitrogen and
          oxygen act as catalysts. S has a cooler core, so it mainly uses the proton--proton
          chain.
  \end{itemize}

  \textbf{(c)} When the hydrogen in the core runs out, fusion in the core stops. The core
  shrinks and heats up, and hydrogen fuses in a shell around it. The outer layers swell and
  cool: the Sun becomes a red giant, like P. When the core reaches about 100 million K,
  helium fuses into carbon by the triple-alpha process ($3\,\pracNuc{4}{2}{He} \to
  \pracNuc{12}{6}{C}$). Some carbon also captures helium to make oxygen. The Sun is not
  massive enough to fuse carbon. Its outer layers drift away as a planetary nebula. The hot
  carbon--oxygen core is left as a white dwarf, like R. It has no fusion and slowly cools.
  \end{modelanswer}
  % layout QA (book build): one extra line lets these criteria fit under the model answer
  % (otherwise they fill only 40% of an extra page).
  \enlargethispage{\baselineskip}
  \begin{markingcriteria}
    \criterion{(a) R is the white dwarf and P is the red giant, each justified with its
      temperature and luminosity (and so its size)}{2}
    \criterion{(a) Both identified without justification, OR one identified and
      justified}{1}
    \criterion{(b) Correct comparison of mass, time on the main sequence AND fusion process
      (Q: more massive, shorter life, CNO cycle; S: less massive, longer life, p--p
      chain)}{3}
    \criterion{(b) Two of these correct}{2}
    \criterion{(b) One of these correct}{1}
    \criterion{(c) Describes all three stages: core hydrogen runs out, so the core contracts
      and heats (hydrogen fuses in a shell; red giant); helium fuses into carbon (and
      oxygen) at about 100 million K; fusion stops, the outer layers drift away as a
      planetary nebula and the core is left as a white dwarf}{3}
    \criterion{(c) Describes two of these stages correctly}{2}
    \criterion{(c) Describes one correct stage}{1}
  \end{markingcriteria}
\end{practiceans}

\end{practiceanswers}
